\documentclass[11pt,letterpaper]{article}
\usepackage[T1]{fontenc}
\usepackage[utf8]{inputenc}
\usepackage{lmodern}
\usepackage[margin=1in,headheight=15pt]{geometry}
\usepackage{amsmath,amssymb,amsthm,mathtools,mathrsfs}
\usepackage{microtype,booktabs,longtable,array,enumitem}
\usepackage[dvipsnames]{xcolor}
\usepackage{fancyhdr}
\usepackage[colorlinks=true,linkcolor=MidnightBlue,citecolor=MidnightBlue,urlcolor=MidnightBlue]{hyperref}
\usepackage[nameinlink,noabbrev]{cleveref}
\hypersetup{pdftitle={Jordan Separation, Cauchy Theory, and Stokes Formulae on the Surcomplex Plane},pdfauthor={Research exposition},pdfsubject={Surcomplex contour integration, definable topology, Hahn differential forms, and residues}}
\pagestyle{fancy}\fancyhf{}
\fancyhead[L]{\small Surcomplex contour calculus}
\fancyhead[R]{\small Separation, periods, and Stokes}
\fancyfoot[C]{\thepage}
\setlength{\parskip}{0.3em}
\setlength{\parindent}{1.3em}
\setlist{nosep,leftmargin=2em}
\allowdisplaybreaks[2]
\emergencystretch=1.5em
\newtheorem{theorem}{Theorem}[section]
\newtheorem{proposition}[theorem]{Proposition}
\newtheorem{lemma}[theorem]{Lemma}
\newtheorem{corollary}[theorem]{Corollary}
\theoremstyle{definition}
\newtheorem{definition}[theorem]{Definition}
\newtheorem{example}[theorem]{Example}
\theoremstyle{remark}
\newtheorem{remark}[theorem]{Remark}
\newcommand{\No}{\mathbf{No}}
\newcommand{\SC}{\mathbf{SC}}
\newcommand{\R}{\mathbb R}
\newcommand{\C}{\mathbb C}
\newcommand{\Q}{\mathbb Q}
\newcommand{\N}{\mathbb N}
\newcommand{\Z}{\mathbb Z}
\newcommand{\OO}{\mathcal O}
\newcommand{\HH}{\mathcal H}
\newcommand{\Acal}{\mathscr A}
\newcommand{\m}{\mathfrak m}
\newcommand{\st}{\operatorname{st}}
\newcommand{\supp}{\operatorname{supp}}
\newcommand{\Res}{\operatorname{Res}}
\newcommand{\res}{\operatorname{res}}
\newcommand{\Ind}{\operatorname{Ind}}
\newcommand{\Tr}{\operatorname{Tr}}
\newcommand{\ord}{\operatorname{ord}}
\newcommand{\id}{\operatorname{id}}
\newcommand{\vH}{v_{\!H}}
\newcommand{\HInt}{\mathop{\int^{\!H}}\nolimits}
\newcommand{\HCoint}{\mathop{\oint^{\!H}}\nolimits}
\newcommand{\sh}{^{\sharp}}
\newcommand{\abs}[1]{\left|#1\right|}
\newcommand{\dR}{\mathrm{dR}}
\newcolumntype{L}[1]{>{\raggedright\arraybackslash}p{#1}}
\numberwithin{equation}{section}
\setcounter{tocdepth}{2}
\begin{document}
\begin{titlepage}
\centering
\vspace*{0.6cm}
{\Large\scshape A research continuation\par}
\vspace{0.65cm}
{\Huge\bfseries Jordan Separation,\par Cauchy Theory,\par and Stokes Formulae\par}
\vspace{0.35cm}
{\LARGE on the Surcomplex Plane\par}
\vspace{0.55cm}
{\Large A Category-Sensitive Contour Calculus\par}
\vspace{0.6cm}
{\large September 21, 2026\par}
\vspace{0.45cm}
\begin{minipage}{0.94\textwidth}
\begin{abstract}
There is no single surcomplex analogue of the Jordan--Cauchy--Stokes package independent of the chosen topology, function algebra, and class of integration chains. The full fine topology admits no nonconstant continuous ordinary real-parameter paths. Nevertheless, ordinary shadow topology and definable topology give two precise separation theorems; Hahn-coherent differential forms give an exact generalized Stokes theorem; and positive-support deformations give nonstandard contours with integer winding and homotopy-invariant periods.

We extend the supplied manuscripts in a more local direction as well. An infinitesimal rescaling converts radius-free Hahn analytic germs into common-domain polynomial-coefficient families. This yields microscopic Cauchy and residue formulae. Combining the supplied finite-algebra deformation theorem with classical residue transformation, we realize its multidimensional residues by explicit separating product tori, even when the original coefficient germs share no ordinary neighborhood. Rational differential forms additionally admit an algebraic period theory on semialgebraic cycles at arbitrary surreal scales. We compare these constructions with definable, Berkovich, perfectoid, $p$-adic, and resurgent integration, keeping their different hypotheses and target objects explicit.
\end{abstract}
\end{minipage}
\vfill
\begin{minipage}{0.94\textwidth}\small
\textbf{Provenance and status.} This article develops and compares constructions rather than asserting an established universal definition of surcomplex integration. Results imported from the two supplied manuscripts are identified as such; published antecedents are cited separately. The separating-torus theorem depends on the supplied finite-algebra theorem and on classical residue transformation. No exhaustive priority claim or formal verification of the general proofs is made.
\end{minipage}
\end{titlepage}
\tableofcontents
\newpage

\section{The question has three independent choices}\label{sec:intro}
The classical Jordan curve theorem concerns separation by an embedded circle. Cauchy's integral theorem concerns periods of holomorphic one-forms. Generalized Stokes concerns the interaction between the boundary of a chain and the exterior derivative of a differential form. These statements cooperate in ordinary complex analysis, but they are not logically interchangeable. In particular, Cauchy's theorem is naturally a homological statement and does not require every contour to be a Jordan curve.

For the surcomplex class field
\[
\SC=\No[i],
\]
three choices must be stated before any counterpart has a definite meaning: the topology in which a curve or domain is considered, the regularity algebra of the integrand, and the definition of the chain and its integral. Real closedness of $\No$ supplies algebraic closedness of $\SC$, but does not select these additional structures.

The supplied manuscript \emph{Surcomplex Analysis} \cite{SA} already distinguishes fine-local analyticity from Hahn-coherent holomorphy and proves a coefficientwise Cauchy theory. The supplied \emph{Infinitesimal Analytic Geometry} \cite{AG} replaces common-domain coefficient functions by coefficient germs with no shared radius and proves a finite-algebra residue theory. The latter explicitly defines residues algebraically rather than assuming an available contour. We retain these distinctions. Our principal extensions are a differential-form Stokes package, a calculus of admissible deformed chains, and a contour realization of those radius-free finite-cluster residues.

\subsection{A map of the answers}
\begin{center}\small
\begin{tabular}{@{}L{0.23\textwidth}L{0.32\textwidth}L{0.35\textwidth}@{}}
\toprule
Framework & Separation or cycles & Integration theorem\\
\midrule
Full fine topology & No nonconstant ordinary continuous paths & No unrestricted fundamental theorem for all fine-analytic functions\\[0.45em]
Standard-part shadow & A thickened ordinary Jordan curve separates into two shadow components & Coherent periods are governed by the ordinary base\\[0.45em]
Definable topology over a real closed field & Definable Jordan separation; field-parameter paths & Algebraic winding and rational periods; additional choices for general integration\\[0.45em]
Hahn differential forms & Ordinary parameter chains, with nonstandard coefficients & Exact generalized Stokes and coefficientwise de Rham theory\\[0.45em]
Positive-support deformations & Infinitesimally moving parameterized chains & Pullback Stokes, homotopy invariance, integer winding\\[0.45em]
Radius-free Hahn germs & Infinitesimal circles and separating product tori & Cauchy, moving-pole residues, and finite-cluster residue realization\\
\bottomrule
\end{tabular}
\end{center}
The affirmative rows are mathematical theorems in specified categories, not approximations to a nonexistent unrestricted fine-topological contour theory.

\subsection{Which assertions come from which source?}
The normal-form and Hahn-field foundations are established mathematics \cite{Poonen,Neumann}. Definable separation is part of real algebraic and o-minimal topology \cite{BO}. Ordinary Stokes and residue transformation are classical inputs \cite{GH,CDS}. Common-domain Cauchy theory, fine-topological obstructions, and moving-pole preparation are developed in \cite{SA}. The radius-free ring, its evaluation and translation, and its finite complete-intersection deformation theorem are developed in \cite{AG}. The deductions in this article are proved using those identified inputs. This is not an independent audit of every theorem in either supplied manuscript.

\section{Workspaces, supports, and the two topological obstructions}\label{sec:foundations}
\subsection{Set-sized fields and strong sums}
Fix a nonzero divisible set-sized ordered subgroup $\Gamma\subset\No$ and put
\begin{equation}\label{eq:workspace}
F=\R((t^\Gamma)),\qquad K=\C((t^\Gamma))=F[i],\qquad
 t^\gamma=\omega^{-\gamma}.
\end{equation}
The real normal-form ordering makes $F$ real closed, and $K$ is algebraically closed. The latter is the Hahn-field algebraic-closure theorem; see \cite[Corollary 4]{Poonen}. An element is a series with well-ordered exponent support. Its valuation is the least exponent, with $v(0)=+\infty$. Define
\[
K^\circ=\{a:v(a)\ge0\},\qquad \m_K=\{a:v(a)>0\}.
\]
Reduction $K^\circ\to\C$ extracts the exponent-zero coefficient. It is the standard-part map $\st$, and every finite element is uniquely $c+\epsilon$, with $c\in\C$ and $\epsilon\in\m_K$.

Full-class neighborhood statements may be interpreted in a class theory such as NBG; the rings, supports, and parameter manifolds used in the integration constructions are set-sized. Every finite collection of surcomplex data lies in a workspace of this kind after adjoining its normal-form supports and taking the divisible hull. Enlarging a workspace is allowed, but all sums and all algebraic constructions below remain set-indexed. Differentiation in a variable $z$ keeps the elements of $K$ fixed. It is not a derivation on $\No$ considered as a transseries field.

\begin{definition}[Strong summability]
A family of Hahn series is strongly summable if the union of its supports is well ordered and only finitely many members contribute to each exponent. Its sum is then defined by finite addition at every exponent. For a well-ordered $E\subset\Gamma_{>0}$, let $E^*$ be its additive monoid, including zero.
\end{definition}
The Hahn--Neumann lemma says that $S+T$ is well ordered for well-ordered $S,T$, with finitely many decompositions of each exponent. Also $E^*$ is well ordered, and only finitely many finite words in $E$ have a prescribed sum. Consequently an expression controlled by $S+E^*$ is legitimate when its contributions are indexed by one source exponent and such a word \cite{Neumann,Poonen}. We use both conclusions, not merely a lower bound on the valuations.

For example, if $\Gamma=\Q+\Q\omega$, the strong sum $\sum_{n\ge0}t^n=(1-t)^{-1}$ exists. Its finite partial sums do not converge to that value in the valuation topology: a tail beginning at exponent $n$ never reaches valuation greater than $\omega$. A fixed positive improvement per iteration need not be cofinal in a higher-rank ordered group. This distinction, emphasized in \cite{AG}, is essential below.

\subsection{Full fine topology does not support ordinary contour paths}
The full fine topology on $\SC$ uses all balls $B(a,r)$, $r\in\No_{>0}$. The following obstruction is from \cite[\S3]{SA}; its short proof is worth retaining.

\begin{proposition}[Set-sized discreteness]\label{prop:fine}
Every set-sized subset of $\SC$ is closed and discrete in the full fine topology. Every fine-continuous map from an ordinary connected real interval into $\SC$ is constant.
\end{proposition}
\begin{proof}
For a point $p$ and a set $A$, the nonzero distances $|p-a|$, $a\in A$, form a set of positive surreals. There is a positive surreal strictly smaller than every member: the cut $\{0\mid |p-a|:a\in A,\ a\ne p\}$ supplies one. Its ball about $p$ meets $A$ only possibly at $p$. Thus $A$ is closed and discrete. The image of an ordinary interval is a set, and a connected subset of a discrete space has at most one point.
\end{proof}
For the same reason, every convergent set-indexed net in this full topology is eventually equal to its limit. The usual real-parameter unit-circle map is not fine-continuous. Calling it a contour below will therefore mean an admissible parameterized integration object, not a fine-continuous loop.

There is a related but distinct obstruction inside a non-Archimedean real closed field with its \emph{intrinsic} order topology. Scaled infinitesimal cosets are clopen and separate points. If $a\ne b$, the coset $a+|b-a|\m_K$ contains $a$ but not $b$. Hence ordinary connectedness is again too weak to model the intended planar geometry. These cosets need not be definable in the language under consideration, which is precisely why definable connectedness can still be useful.

\begin{remark}[Three topologies must not be conflated]
The intrinsic topology of $K$ uses radii in $F$. The topology inherited by the set $K$ from the full class $\SC$ also permits radii outside $F$ and is discrete by \cref{prop:fine}. The standard-part shadow topology introduced below is coarser still. A claim about convergence or connectedness in one of them is not automatically a claim in either of the others.
\end{remark}

\subsection{An unrestricted fundamental theorem is impossible}
\begin{proposition}[A no-go theorem for endpoint integration]\label{prop:noFTC}
There is no family of linear endpoint integral functionals, defined for all fine-locally analytic derivatives and arbitrary endpoints, satisfying
\[
I_{a,b}(H')=H(b)-H(a)
\]
for every fine-locally analytic function $H$ on the full surcomplex plane.
\end{proposition}
\begin{proof}
Let $q(z)=0$ on the infinitesimals and $q(z)=1$ elsewhere. The infinitesimals and all their additive cosets are fine-open, so $q$ is locally constant and hence fine-analytic, with $q'=0$. At $a=0$, $b=1$, the proposed identity gives $I_{0,1}(0)=1$, contradicting linearity.
\end{proof}
The supplied manuscript also constructs a full fine-analytic logarithm with derivative $1/z$, while its coherent contour period around the ordinary unit circle is $2\pi i$ \cite[\S\S6,11]{SA}. This is not a contradiction: that logarithm is not a coherent primitive on the punctured ordinary base. A valid fundamental theorem must restrict the primitive and the integration category together.

\section{Two precise Jordan counterparts}\label{sec:jordan}
\subsection{A shadow Jordan theorem}
For ordinary $U\subset\C$, write $U\sh=\st^{-1}(U)\subset K^\circ$. Give $K^\circ$ the topology whose open sets are precisely $\st^{-1}(V)$ for ordinary open $V\subset\C$. We call this the \emph{shadow topology}. Its points within a standard-part fiber are topologically indistinguishable; its natural quotient is the ordinary complex plane.

\begin{theorem}[Jordan separation in the shadow topology]\label{thm:shadow}
Let $J\subset\C$ be an ordinary Jordan curve, with bounded interior $D$ and exterior $E$. Then
\[
K^\circ\setminus J\sh=D\sh\sqcup E\sh
\]
has exactly two connected components in the shadow topology, and both have boundary $J\sh$.
\end{theorem}
\begin{proof}
Every open set is saturated for the standard-part fibers. A separation of $D\sh$ would therefore descend to a separation of $D$, and similarly for $E$. The ordinary Jordan theorem makes these two sets connected. They are disjoint open subsets of the complement and cover it, so they are its two components. Neighborhoods are inverse images of ordinary neighborhoods, and every neighborhood of a point of $J$ meets both $D$ and $E$. This proves the boundary assertions.
\end{proof}
The separating object is $J\sh$, a \emph{thickened} Jordan curve, not the set of ordinary points $J$ alone. Neither component is being called fine-connected. Also ``exterior'' here means unbounded in the ordinary shadow: all of $K^\circ$ is bounded by any fixed positive infinite element of $F$. The same construction transports to an affine workspace $a+rK^\circ$.

This theorem is a quotient-level analogue. It is particularly appropriate for the coherent functions of \cite{SA}, which are specified over an ordinary base and evaluated on its infinitesimal thickening. It cannot distinguish two zeros lying in the same standard-part fiber; a finer scale is needed for that.

\subsection{Definable Jordan separation over a real closed field}
The second counterpart retains the actual field points but changes connectedness. A semialgebraic set is defined using finitely many polynomial equalities and inequalities with coefficients in $F$. A set is \emph{definably connected} if it has no separation into two nonempty relatively open definable subsets. Paths are continuous definable maps from $[0,1]_F$, not maps from the ordinary real interval.

\begin{theorem}[Semialgebraic Jordan theorem; established counterpart]\label{thm:defJordan}
Let $F$ be a real closed field and let $J\subset F^2$ be a semialgebraic embedded circle. Its complement has exactly two semialgebraically connected components, one bounded and one unbounded, and $J$ is the common boundary. More generally, smooth definable Jordan--Brouwer separation holds in the o-minimal manifold framework of Berarducci--Otero \cite{BO}.
\end{theorem}
For clarity, the generalization in the last sentence refers to the differentiability and definability hypotheses of that framework; it is not a statement about arbitrary class functions or arbitrary fine-topological subsets.

One can understand the semialgebraic assertion through finite compatible triangulation. Compactify the plane by a semialgebraic sphere, triangulate the embedded circle together with the ambient sphere, and apply planar combinatorial separation to the resulting finite pair. Its finite incidence data do not depend on the Archimedean property of the coefficient field. In the polygonal case, a generic ray-crossing test supplies the same two sides, and finite polygonal paths establish their definable connectedness. The o-minimal intersection-theoretic approach instead constructs degree and separation directly \cite{BO}; semialgebraic triangulation and its differentiable refinements are discussed in \cite{OS}.

Circles with infinitesimal radius and polygons with infinite vertices are included. Thus this is a genuine geometric theorem over the non-Archimedean field, not merely a theorem about standard parts. It does not contradict \cref{prop:fine}: $[0,1]_F$ is definably connected, but need not be connected as an ordinary topological space.

\subsection{What Jordan separation does not supply}
Separation by itself does not define an integral. Even in a definable setting, integrating a definable function can require new functions or a larger target object. Conversely, a useful period theory can be defined on cycles with self-intersections, without invoking Jordan separation. We will use the shadow theorem to interpret ordinary-base domains, the definable theorem to interpret algebraic winding, and parameterized chains to prove Stokes. These are compatible uses, but not identifications of categories.

\section{The inherited Hahn-coherent Cauchy theory}\label{sec:cauchy}
\subsection{Common-domain coefficient functions}
Let $U\subset\C$ be ordinary and open. In a fixed workspace put
\[
\HH(U)=\OO(U)((t^\Gamma)).
\]
An element $G=\sum_{\lambda\in S}g_\lambda t^\lambda$ has a single well-ordered support and ordinary holomorphic coefficients on the same $U$. For $z=c+\epsilon\in U\sh$, define
\begin{equation}\label{eq:evaluate}
G\sh(c+\epsilon)=\sum_{\lambda,n\ge0}
 t^\lambda\frac{g_\lambda^{(n)}(c)}{n!}\epsilon^n.
\end{equation}
The controlling support is $S+(\supp\epsilon)^*$. Differentiation acts on the $g_\lambda$. These definitions and their injectivity, composition, and gluing properties are supplied in \cite[\S5]{SA}.

For an ordinary piecewise $C^1$ path $\gamma$ in $U$, define
\begin{equation}\label{eq:Hintegral}
\HInt_\gamma G(z)\,dz
 =\sum_\lambda t^\lambda\int_\gamma g_\lambda(z)\,dz.
\end{equation}
The support cannot increase. This is $K$-linear: multiplication by a fixed Hahn constant gives finite convolution at each exponent, so it commutes with coefficientwise integration. Concatenation, orientation reversal, and ordinary changes of parameter also hold coefficientwise.

\begin{theorem}[Coherent Cauchy theorem and formula; from \cite{SA}]\label{thm:Cauchy}
If an ordinary closed cycle $\gamma$ is null-homologous in $U$ and $G\in\HH(U)$, then $\HCoint_\gamma G\,dz=0$. If $D$ is a bounded ordinary Jordan domain with piecewise $C^1$ boundary and $\overline D\subset U$, then for $z\in D\sh$,
\begin{equation}\label{eq:Cauchy}
(D^kG)\sh(z)=\frac{k!}{2\pi i}\HCoint_{\partial D}
 \frac{G(\zeta)}{(\zeta-z)^{k+1}}\,d\zeta,
 \qquad k\ge0.
\end{equation}
The boundary is positively oriented, and the integrand is expanded as Hahn data near that ordinary boundary.
\end{theorem}
\begin{proof}
Cauchy's theorem applies to every coefficient. For the formula, write $z=c+\epsilon$, and expand
\[
(\zeta-c-\epsilon)^{-k-1}
=\sum_{j\ge0}\binom{k+j}{j}
 \frac{\epsilon^j}{(\zeta-c)^{k+j+1}}.
\]
The base denominator is bounded away from zero on the ordinary contour. Strong summability allows finite regrouping at every Hahn exponent. Ordinary Cauchy formulae produce precisely the Taylor evaluation of $D^kG$ at $c+\epsilon$.
\end{proof}
Null-homology is a real hypothesis. On $\C^\times$, the coherent form $dz/z$ has unit-circle period $2\pi i$, so ``closed contour'' alone does not imply a zero integral.

\subsection{Primitives and displaced endpoints}
On a simply connected ordinary $U$, integrating each coefficient with a fixed normalization gives a coherent primitive. On a general connected $U$, a coherent primitive exists exactly when all ordinary closed-cycle Hahn periods vanish. All coherent primitives of the same function differ by one constant in $K$ \cite[\S6]{SA}.

For finite nonstandard endpoints $z_j=a_j+\epsilon_j$ and an ordinary path $\gamma:a_1\to a_2$, the endpoint correction is
\[
T_G(a,\epsilon)=\sum_{n\ge0}\frac{D^nG(a)}{(n+1)!}\epsilon^{n+1}.
\]
The integral
\begin{equation}\label{eq:endpoint}
-T_G(a_1,\epsilon_1)+\HInt_\gamma G\,dz+T_G(a_2,\epsilon_2)
\end{equation}
equals the difference of a coherent primitive at the endpoints. It gives a useful fundamental theorem without contradicting \cref{prop:noFTC}, because the primitive class has been restricted.

\subsection{Order-valued estimates}
The coefficientwise definition does not discard the surreal order. A continuous real Hahn family that is pointwise nonnegative on an ordinary compact interval has nonnegative Hahn integral: its first nonzero coefficient function is nonnegative and has positive ordinary integral. Rotating a nonzero complex integral by the constant $\overline I/|I|$ then gives the coherent ML inequality
\[
\left|\HInt_\gamma G\,dz\right|\le L M
\]
when the ordinary length is $L$ and $|G|\le M\in F_{\ge0}$ along the parameterized contour. The full argument and the resulting nonstandard Cauchy estimates appear in \cite{SA}. This is an ordered-field estimate, not a claim that fine Riemann sums converge.

\section{A generalized Hahn--Stokes theorem}\label{sec:Stokes}
\subsection{The differential-form complex}
Let $M$ be an ordinary smooth manifold. Define the complex of globally supported Hahn forms by
\begin{equation}\label{eq:Hforms}
\Omega_H^q(M)=\Omega^q(M;\C)((t^\Gamma)),\qquad
 d\left(\sum_\lambda\omega_\lambda t^\lambda\right)
 =\sum_\lambda d\omega_\lambda t^\lambda.
\end{equation}
Wedge products use finite Hahn convolution. Hence $d^2=0$ and the graded Leibniz rule hold. Ordinary smooth pullback commutes with $d$ and does not enlarge support. At this stage all coefficient forms live on an ordinary manifold; evaluating a smooth coefficient at a nonstandard target point has not been assumed.

For a compact oriented ordinary $q$-manifold $C$, or a finite smooth singular $q$-chain, set
\[
\HInt_C\omega=\sum_\lambda t^\lambda\int_C\omega_\lambda.
\]
Each coefficient integral is ordinary. The output lies in $K$, even when no real-valued bound on the entire coefficient family is available.

\begin{theorem}[Generalized Hahn--Stokes]\label{thm:Stokes}
Let $M$ be a compact oriented ordinary smooth $m$-manifold with boundary, and let $\omega\in\Omega_H^{m-1}(M)$ be smooth up to the boundary. With the usual outward-normal-first boundary orientation,
\begin{equation}\label{eq:Stokes}
\boxed{\qquad \HInt_M d\omega=\HInt_{\partial M}\omega.\qquad}
\end{equation}
The same identity holds for finite smooth chains, with their chain boundary.
\end{theorem}
\begin{proof}
For every exponent $\lambda$, ordinary generalized Stokes gives
$\int_Md\omega_\lambda=\int_{\partial M}\omega_\lambda$; see, for example, \cite[Theorem 4.6.1]{GH}. Both Hahn sums have support contained in $\supp\omega$. Equality at each coefficient proves \eqref{eq:Stokes}. On finite chains, the same proof and cancellation of oppositely oriented faces apply.
\end{proof}
This is a full differential-form theorem in every ordinary finite dimension. Its content is not an assertion that the full surcomplex plane is an ordinary manifold. The chain is ordinary; its coefficients and its integral may be genuinely nonstandard. The pullback construction in \cref{sec:deformations} will enlarge the chains themselves.

\subsection{Green and divergence formulae}
On an ordinary planar domain, for a $C^1$ Hahn coefficient family $G$,
\[
d(G\,dz)=2i\frac{\partial G}{\partial\bar z}\,dx\wedge dy.
\]
Thus \cref{thm:Stokes} gives
\begin{equation}\label{eq:Green}
\HCoint_{\partial D}G\,dz
=2i\HInt_D\frac{\partial G}{\partial\bar z}\,dx\,dy.
\end{equation}
For holomorphic coefficients the right side is zero. This supplies a differential-form explanation of the coherent Cauchy theorem for bounding chains.

Likewise, for a Hahn family of vector fields $X$ on an ordinary oriented Euclidean region with volume form $\mathrm{vol}$,
\[
d(\iota_X\mathrm{vol})=(\operatorname{div}X)\mathrm{vol}.
\]
Its boundary flux equals the Hahn integral of its divergence. No separate limiting construction is needed for either formula: both are specializations of the same chain identity.

\subsection{Cohomology and the exact source of periods}
\begin{theorem}[Hahn de Rham cohomology]\label{thm:deRham}
For the globally supported complex \eqref{eq:Hforms},
\begin{equation}\label{eq:deRham}
H^q(\Omega_H^\bullet(M),d)
\simeq H^q_{\dR}(M;\C)((t^\Gamma)).
\end{equation}
If $H^q_{\dR}(M;\C)$ is finite-dimensional, the right side is naturally
$H^q_{\dR}(M;\C)\otimes_\C K$.
\end{theorem}
\begin{proof}
A Hahn form is closed exactly when each coefficient is closed. Send it to the Hahn family of coefficient cohomology classes. For surjectivity, choose a complex-linear representative map from ordinary cohomology to closed forms and apply it coefficientwise. If every coefficient class vanishes, choose a complex-linear right inverse of $d$ on its image. Applied coefficientwise, it produces a primitive without increasing support. This proves the kernel is exactly the exact Hahn forms. A finite basis identifies the Hahn extension of a finite-dimensional vector space with its tensor product with $K$.
\end{proof}
The choices in the proof establish existence; the map on cohomology itself is canonical. In infinite dimension, the tensor-product identification should not be asserted without qualification. Nor does the global-support convention automatically produce an unrestricted sheaf of smooth Hahn families: locally chosen supports can fail to have well-ordered union.

By the ordinary de Rham theorem, a closed Hahn form of positive degree has a global Hahn primitive exactly when all its ordinary cycle periods vanish. On a punctured ordinary plane the class of $dz/z$ survives; on an ordinary contractible domain positive-degree cohomology vanishes. Thus the topological information in coherent periods is inherited from the \emph{ordinary base}, not from fine-connected components of $\SC$.

\section{Infinitesimally deformed chains and exact winding}\label{sec:deformations}
\subsection{Admissible pullbacks}
There is no need to restrict every contour to a fixed ordinary image. The support mechanism permits controlled nonstandard movement of the image itself.

\begin{definition}[Admissible positive-support deformation]\label{def:deformation}
Let $M$ be an ordinary compact smooth parameter manifold and $U\subset\R^d$ an ordinary open set. An admissible deformation of a smooth map $\phi_0:M\to U$ is
\begin{equation}\label{eq:deformation}
\phi=\phi_0+\eta,\qquad
\eta=\sum_{e\in E}\eta_e t^e,
\end{equation}
where $E\subset\Gamma_{>0}$ is well ordered and the $\eta_e$ are smooth real vector-valued coefficient functions on $M$. For complex coordinates use their real and imaginary parts. Piecewise smooth parameterizations are handled using a finite subdivision.
\end{definition}
Let $\omega=\sum_\lambda\omega_\lambda t^\lambda$ have real-analytic coefficient forms on a common target neighborhood $U$. Pull back its coefficient functions by their Taylor series about $\phi_0$, and substitute
\[
d\phi_j=d(\phi_0)_j+\sum_e d(\eta_e)_j t^e
\]
for the coordinate differentials. For a coefficient function $a$, the first operation is
\[
a(\phi_0+\eta)=\sum_{\alpha\in\N^d}
 \frac{(\partial^\alpha a)\circ\phi_0}{\alpha!}\eta^\alpha.
\]
The resulting object is an ordinary smooth Hahn form on $M$.

\begin{lemma}[Support-controlled pullback]\label{lem:pullback}
If $S=\supp\omega$, the pullback just defined has support contained in $S+E^*$, with finitely many contributions at each exponent. It satisfies
\[
d(\phi^*\omega)=\phi^*(d\omega),\qquad
\phi^*(\omega\wedge\nu)=\phi^*\omega\wedge\phi^*\nu.
\]
It agrees with canonical real-analytic infinitesimal evaluation of the coefficients at every ordinary parameter point.
\end{lemma}
\begin{proof}
A term uses a source exponent from $S$ and finitely many factors drawn from $\eta$ and $d\eta$. Its exponent belongs to $S+E^*$. The support lemma gives only finitely many words and hence only finitely many Taylor indices and coefficient products at a prescribed exponent. Each coefficient is therefore a finite sum of ordinary smooth forms. Ordinary Taylor chain and product identities apply coefficientwise; every rearrangement in their proofs is finite at that exponent. This proves the identities and the evaluation assertion.
\end{proof}
Real analyticity is imposed on the \emph{target} coefficient functions; the parameter coefficients need only be smooth. Assigning values of an arbitrary smooth function to infinitesimal arguments is an additional extension problem. Its Taylor jet alone does not specify a unique extension of the original real function. We do not silently make such an extension choice.

\subsection{Stokes and homotopy for moving chains}
Define the integral along the parameterized chain by
\[
\HInt_\phi\omega:=\HInt_M\phi^*\omega.
\]
This is a definition on a parameterized object, not an attempt to integrate over its image using a fine-topological measure.

\begin{theorem}[Stokes for admissible deformations]\label{thm:movingStokes}
For an admissible $\phi$ on a compact oriented $m$-manifold with boundary and an admissible real-analytic Hahn $(m-1)$-form $\omega$,
\begin{equation}\label{eq:movingStokes}
\HInt_\phi d\omega=\HInt_{\phi|_{\partial M}}\omega.
\end{equation}
If $\Phi$ is an admissible homotopy of closed chains and $d\omega=0$ on its target neighborhood, the periods at its two ends are equal.
\end{theorem}
\begin{proof}
Apply \cref{thm:Stokes} to $\phi^*\omega$, using \cref{lem:pullback}. For the homotopy, apply the same identity to the parameter cylinder $[0,1]\times M$; its end faces have opposite orientations. The pullback of $d\omega$ is zero, and the side boundary is absent for a closed chain.
\end{proof}
Meromorphic forms are allowed when their denominators have nonvanishing leading coefficient on the base image. In that case geometric inversion supplies holomorphic Hahn data on a common ordinary neighborhood of the image, so the theorem applies. Passing through a leading pole is not justified by calling its displacement infinitesimal.

In one complex variable, a useful direct identity is
\begin{equation}\label{eq:variation}
\partial_u\bigl(G(\gamma_u)\partial_s\gamma_u\bigr)
=\partial_s\bigl(G(\gamma_u)\partial_u\gamma_u\bigr)
\end{equation}
for a fixed holomorphic Hahn $G$ and an admissible two-parameter contour. These derivatives differentiate the ordinary coefficient functions of the parameter, not a putative fine-continuous real-time path. Integration in $s$ shows that every correction to a closed-contour period is exactly zero. The conclusion is equality in $K$, not equality only after standard part.

Ordinary reparameterization invariance is immediate. Admissible analytic coordinate changes also preserve the integral whenever both pullbacks are defined: the composition identity follows by the same support argument. More generally, two admissible descriptions with the same pointwise pulled-back differential have identical Hahn coefficients by normal-form uniqueness and hence the same integral. None of these assertions grants automatic compatibility to arbitrary unrelated scale atlases.

\subsection{An integer winding theorem at moving points}
\begin{theorem}[Winding survives infinitesimal movement]\label{thm:winding}
Let $\gamma=\gamma_0+\eta$ be an admissible deformation of an ordinary closed piecewise smooth contour, and let $a=a_0+\xi$, $\xi\in\m_K$, with $a_0\notin\gamma_0$. Then
\begin{equation}\label{eq:winding}
\frac1{2\pi i}\HCoint_\gamma\frac{dz}{z-a}
=\Ind_{\C}(\gamma_0,a_0)\in\Z.
\end{equation}
\end{theorem}
\begin{proof}
On the parameter cylinder define
\[
g_u(s)=\gamma_0(s)-a_0+u(\eta(s)-\xi).
\]
Its leading coefficient never vanishes; compactness of the ordinary contour also supplies a positive ordinary separation from $a_0$. Its inverse is therefore a uniformly supported Hahn family. Direct differentiation gives
\[
\partial_u\left(\frac{\partial_sg_u}{g_u}\right)
=\partial_s\left(\frac{\partial_ug_u}{g_u}\right).
\]
The integral of the right side along the closed parameter is zero, coefficient by coefficient. The period at $u=1$ equals that at $u=0$, which is the ordinary winding integral.
\end{proof}
This theorem explains why an integer, rather than an arbitrary surreal number, appears in coherent argument principles. It does not define an index when the standard part of a pole lies on the base contour. Such a situation can sometimes be resolved in a smaller affine chart, but needs a new admissibility check.

\section{Radius-free germs acquire contours after infinitesimal scaling}\label{sec:radiusfree}
\subsection{The local coefficient algebra}
The second supplied manuscript uses
\begin{equation}\label{eq:An}
R_n=\C\{z_1,\ldots,z_n\},\qquad
\Acal_n=R_n((t^\Gamma)).
\end{equation}
Each coefficient is an ordinary convergent germ at zero, but their convergence radii may tend to zero. For example,
\begin{equation}\label{eq:shrinking}
H(z)=\sum_{r\ge1}\frac{t^r}{1-rz}\in\Acal_1
\end{equation}
has no common ordinary disk of holomorphy. It nevertheless evaluates at every $z\in\m_K$. The support-controlled evaluation, translation, division, and finite-algebra theorems for this ring are inputs from \cite{AG}.

Consequently, applying \cref{thm:Cauchy} to \eqref{eq:shrinking} on a fixed ordinary circle would leave a gap. The following rescaling, rather than an imaginary uniform radius, removes it.

\begin{lemma}[Polynomial coefficients after infinitesimal scaling]\label{lem:scaling}
Let $H\in\Acal_n$, $a\in\m_K^n$, and $s_1,\ldots,s_n\in\Gamma_{>0}$. Set $r_i=t^{s_i}$. Then
\begin{equation}\label{eq:scaled}
H(a_1+r_1w_1,\ldots,a_n+r_nw_n)
=\sum_\lambda P_\lambda(w)t^\lambda,
\qquad P_\lambda\in\C[w_1,\ldots,w_n],
\end{equation}
is a Hahn family with a common ordinary domain $\C^n$. It evaluates correctly at every finite tuple $w\in(K^\circ)^n$.
\end{lemma}
\begin{proof}
Translate the germ by the infinitesimal tuple $a$ using the strong Taylor translation of \cite[\S3]{AG}. Expand each translated coefficient in the variables $r_iw_i$. At a fixed exponent, terms are indexed by a source exponent together with a finite word from
\[
E_a=\bigcup_i\supp(a_i)\ \cup\ \{s_1,\ldots,s_n\}.
\]
The support lemma gives finitely many such words. In particular only finitely many Taylor degrees occur at that exponent, so its coefficient is a polynomial in $w$. The support is contained in $\supp H+E_a^*$. Evaluation at a finite $w=c+\epsilon$ is justified by adjoining the positive supports of its infinitesimal parts. Polynomial evaluation and Taylor regrouping give the asserted identity.
\end{proof}
The scalar form of this lemma occurs in \cite[Lemma 10.1]{AG}; the proof above makes the separate-coordinate version and its contour consequence explicit. No radius is claimed to work before rescaling. The resulting coefficients need not have any uniform degree or growth bound as the Hahn exponent varies.

\subsection{Microscopic Cauchy formulae}
For $r=t^s$, the symbol $\{|z|=r\}$ in a Hahn integral will mean the ordinary parameterized circle $z=r e^{i\theta}$, $0\le\theta\le2\pi$, with its pulled-back differential. It does not mean the full class of points of that modulus, equipped with a fine-topological integration theory.

\begin{theorem}[Cauchy formula for a radius-free germ]\label{thm:localCauchy}
Let $H\in\Acal_1$ and $a\in\m_K$. Choose $s>0$ with $s<v(a)$ when $a\ne0$, and let $r=t^s$. For every $k\ge0$,
\begin{equation}\label{eq:localCauchy}
\boxed{\quad H^{(k)}(a)=\frac{k!}{2\pi i}
\HCoint_{|z|=r}\frac{H(z)}{(z-a)^{k+1}}\,dz.\quad}
\end{equation}
If $a=0$, any $s\in\Gamma_{>0}$ is allowed.
\end{theorem}
\begin{proof}
By \cref{lem:scaling}, $G(w)=H(rw)$ has ordinary entire polynomial coefficients. The normalized point $a/r$ is infinitesimal. Apply \cref{thm:Cauchy} to $G$ on the ordinary unit circle. Since $G^{(k)}(a/r)=r^kH^{(k)}(a)$ and $dz=r\,dw$, the scaling factors yield \eqref{eq:localCauchy}.
\end{proof}
The same proof gives a product-torus Cauchy formula for $H\in\Acal_n$, using one sufficiently large infinitesimal circle around each coordinate of the evaluation point. A product of $n$ circles is a real $n$-dimensional cycle, often called the distinguished boundary of a polydisk. It is \emph{not} the whole real $(2n-1)$-dimensional boundary of that polydisk.

There is also a local primitive theorem. Applying ordinary germwise integration coefficient by coefficient makes $D:\Acal_1\to\Acal_1$ surjective, with kernel $K$. More generally, the ordinary analytic Poincare homotopy applied coefficientwise makes the analytic differential-form complex of $\Acal_n$ exact in positive degrees. Thus microscopic periods of nonsingular forms vanish, while poles supply residue obstructions. These assertions use the prescribed Hahn-germ algebra, not all fine-analytic functions.

\subsection{Choosing a separating microscopic circle}
\begin{lemma}[A radius protecting an entire one-variable cluster]\label{lem:protect}
Suppose $G=g+E\in\Acal_1$, where $g$ has ordinary order $m\ge1$ at zero and $\delta=\vH(E)>0$. Choose $s>0$ with $ms<\delta$, and put $r=t^s$. Then $r^{-m}G(rw)$ has leading coefficient $u(0)w^m$, where $g(z)=z^mu(z)$ and $u(0)\ne0$. It has no zeros on the ordinary unit-circle parameter. Every infinitesimal zero $a$ of $G$ satisfies $v(a)>s$.
\end{lemma}
\begin{proof}
The higher powers in $u(rw)-u(0)$ have positive valuation. Every term from $r^{-m}E(rw)$ has valuation at least $\delta-ms>0$. This proves the leading-coefficient and boundary assertions. For a nonzero infinitesimal zero with $\rho=v(a)\le s$, the term $g(a)$ has valuation $m\rho\le ms<\delta$, while $E(a)$ has valuation at least $\delta$. They cannot cancel. The zero $a=0$ satisfies the conclusion by convention.
\end{proof}
When $E=0$, every positive $s$ works. Divisibility of $\Gamma$ supplies the strictly smaller positive $s$ in the nonzero case, even when $\Gamma$ has arbitrary rank.

\begin{theorem}[Radius-free moving-pole residue theorem]\label{thm:localresidue}
Let $N,G\in\Acal_1$, with $G\ne0$. Normalize its leading exponent and choose a protecting microscopic circle as in \cref{lem:protect} when its leading germ vanishes at zero. If the leading germ is a unit, choose any sufficiently small infinitesimal circle. Then
\begin{equation}\label{eq:localresidue}
\frac1{2\pi i}\HCoint_{|z|=r}\frac{N(z)}{G(z)}\,dz
=\sum_{a\in\m_K}\res_a\left(\frac NG\,dz\right),
\end{equation}
where the right side is the finite sum over actual poles. Every pole of the quotient in the monad lies inside the chosen normalized circle.
\end{theorem}
\begin{proof}
If the leading germ is a unit, then $G$ is a unit of $\Acal_1$ and the assertion follows from \cref{thm:localCauchy} or coefficientwise primitives. Otherwise Hahn-germ preparation and division from \cite[\S4]{AG} give
\[
\frac NG=J+\frac RP,\qquad J\in\Acal_1,\quad
\deg R<\deg P=m,
\]
where $P$ is monic and all its roots are infinitesimal. The normalization absorbs the leading monomial and unit factors into the numerator. These roots are the zeros of $G$, with multiplicities, and satisfy the protecting bound. Factor $P$ in the algebraically closed field $K$ and use finite partial fractions
\[
\frac RP=\sum_a\sum_{j=1}^{m_a}\frac{c_{a,j}}{(z-a)^j}.
\]
The $J$ integral is zero. For each $a$, $a/r$ is infinitesimal, so expansion on $z=rw$, $|w|=1$, is strongly summable. Its normalized integral is $1$ for $(z-a)^{-1}dz$ and zero for every $(z-a)^{-j}dz$, $j>1$. Thus the integral is $\sum_a c_{a,1}$, exactly the sum of local residues, with removable points contributing zero.
\end{proof}
This proof extends the common-domain moving-pole theorem of \cite{SA} to the radius-free setting. It does not require a common ordinary representative for the original infinitely many coefficient germs.

The corresponding algebraic expansion is also useful. Let $\mathcal M_0$ be the field of ordinary meromorphic germs at zero. For normalized $G=g+E$, the quotient has an expansion in $\mathcal M_0((t^\Gamma))$, and its coefficientwise residue is
\begin{equation}\label{eq:rho}
\rho(N/G)=\sum_{k\ge0}(-1)^k
 \res_0\left(\frac{N E^k}{g^{k+1}}\,dz\right).
\end{equation}
It is strongly summable and agrees with \eqref{eq:localresidue}; preparation and the same partial fractions establish the agreement. An arbitrary Hahn family of meromorphic germs is a larger object than a quotient of two elements of $\Acal_1$. We do not assert that every such arbitrary family admits the microscopic contour construction.

\section{Separating tori for finite analytic zero schemes}\label{sec:tori}
This section addresses the most specific gap between the two supplied manuscripts: the several-variable residue in \cite{AG} is algebraically defined, without asserting a common contour for all its coefficient germs. We now construct a contour representation after changing to separating denominators. The denominator change is essential.

\subsection{The finite-algebra input}
Assume
\begin{equation}\label{eq:CI}
F_i=f_i+E_i\in\Acal_n,\qquad \vH(E_i)>0,
\qquad
\mu=\dim_\C\frac{R_n}{(f_1,\ldots,f_n)}<\infty,
\end{equation}
where $f_i(0)=0$ and $\mu>0$. Zero perturbations are permitted. This is an isolated ordinary complete intersection perturbed by positive Hahn support. Let $E$ be the union of the nonzero perturbation supports.

We import the following package proved in the supplied manuscript \cite[\S\S7--9]{AG}. The quotient
\[
B_F=\Acal_n/(F_1,\ldots,F_n)
\]
has $K$-dimension $\mu$. Any chosen ordinary monomial basis remains a basis, and its multiplication matrices and normal-form operations have explicit support in $E^*$. A residue functional is defined by
\begin{equation}\label{eq:AGresidue}
\Res_F(H)=\sum_{k\in\N^n}(-1)^{|k|}
 \res_{f^{k+1}}(H E_1^{k_1}\cdots E_n^{k_n}),
\end{equation}
where the ordinary residue includes the normalization $(2\pi i)^{-n}$ and
$f^{k+1}=(f_1^{k_1+1},\ldots,f_n^{k_n+1})$. It is strongly summable, has support contained in $\supp H+E^*$, and annihilates $(F)$. In particular,
\begin{equation}\label{eq:traceinput}
\Res_F(HJ_F)=\Tr(M_H)=\sum_a e_aH(a),\qquad
J_F=\det\left(\frac{\partial F_i}{\partial z_j}\right),\quad
\sum_a e_a=\mu.
\end{equation}
The sums over points are finite. These results are inputs, not independently reproved here. Our contour theorem is a consequence of them, the explicit support-controlled normal form, and classical residue transformation.

\subsection{Univariate separating denominators}
Let $M_i$ be the multiplication matrix of $z_i$ in the fixed basis of $B_F$. Define
\begin{equation}\label{eq:Qi}
Q_i(T)=\det(TI-M_i)
=T^\mu+\sum_{j<\mu}q_{ij}T^j.
\end{equation}
Its exponent-zero matrix is multiplication by $z_i$ in the ordinary local algebra $R_n/(f)$, hence is nilpotent. Therefore every nonzero $q_{ij}$ has positive valuation.

By Cayley--Hamilton, $Q_i(z_i)$ belongs to $(F)$. Apply the supplied support-controlled normal form to this element. Since its remainder is zero, it provides a matrix $\mathsf C=(C_{ij})$ with
\begin{equation}\label{eq:matrixC}
Q_i(z_i)=\sum_{j=1}^n C_{ij}(z)F_j(z),\qquad
C_{ij}\in\Acal_n,
\end{equation}
and all entries have nonnegative support contained in $E^*$. This last support statement is useful: membership in the ideal alone would not have exhibited an appropriate coefficient bound for an arbitrary choice of representation.

Choose $s_i\in\Gamma_{>0}$ satisfying the finitely many inequalities
\begin{equation}\label{eq:torusradius}
(\mu-j)s_i<v(q_{ij})\quad
\text{for every nonzero }q_{ij},\qquad r_i=t^{s_i}.
\end{equation}
Such a choice exists by divisibility. If all lower coefficients of a particular $Q_i$ vanish, any positive $s_i$ works. The scaled denominator is
\[
r_i^{-\mu}Q_i(r_iw_i)
=w_i^\mu+\sum_{j<\mu}q_{ij}r_i^{j-\mu}w_i^j,
\]
with strictly positive valuations in every non-leading coefficient. It is therefore invertible as common-domain Hahn data near the ordinary unit circle.
Every root $b$ of $Q_i$ satisfies $v(b)>s_i$: otherwise the leading term $b^\mu$ has strictly smaller valuation than each $q_{ij}b^j$, by \eqref{eq:torusradius}, and cannot cancel. Thus every common zero of $F$ lies strictly inside the separating coordinate circles. The additional points in the product zero set of the $Q_i$ are handled by the residue-transformation determinant, not counted as additional zeros of $F$.

\subsection{The contour realization theorem}
\begin{theorem}[Separating-torus realization]\label{thm:torus}
Assume the finite-algebra and normal-form results just stated from \cite{AG}. Let $Q_i$, $\mathsf C$, and $r_i$ be as in \eqref{eq:Qi}--\eqref{eq:torusradius}. For every $H\in\Acal_n$,
\begin{equation}\label{eq:torus}
\boxed{\quad
\Res_F(H)=\frac1{(2\pi i)^n}
\HInt_{\mathbb T_r}
 \frac{H(z)\det\mathsf C(z)}{\prod_{i=1}^nQ_i(z_i)}
 \,dz_1\wedge\cdots\wedge dz_n,
\quad}
\end{equation}
where $\mathbb T_r$ is parameterized by $z_i=r_i e^{i\theta_i}$ and oriented by $d\theta_1\wedge\cdots\wedge d\theta_n$. The integrand has a well-defined Hahn pullback with ordinary smooth coefficient forms. The result is independent of the allowed radii and of the chosen representing matrix in \eqref{eq:matrixC}. Its final support is contained in $\supp H+E^*$.
\end{theorem}
\begin{proof}
There are three points to verify.

\emph{Existence of the integral.} By \cref{lem:scaling}, $H(rw)$ and every entry of $\mathsf C(rw)$ have polynomial complex coefficients at each Hahn exponent. Condition \eqref{eq:torusradius} makes each reciprocal denominator a strong geometric expansion on a neighborhood of the ordinary product unit torus. Products and the differential factors $r_i\,dw_i$ remain strongly supported. Thus the pullback is an ordinary smooth Hahn $n$-form on a compact real $n$-torus, and its coefficientwise integral exists.

\emph{Residue transformation.} The classical local transformation rule is
\begin{equation}\label{eq:transform}
\Res_F(H)=\Res_Q(H\det\mathsf C)
\quad\text{when }Q=\mathsf C F.
\end{equation}
Its ordinary form, including the determinant factor, is standard residue theory \cite[\S0, equation (0.6)]{CDS}. To justify its use here, fix a Hahn exponent. The normal-form, determinant, and residue constructions depend at that exponent on finitely many ordinary coefficient germs and finitely many positive-support words. Replace those finitely many perturbation terms by independent ordinary parameters. All operations then have formal parameter expansions, and the identity $Q=\mathsf C F$ holds formally. Truncate at an order higher than every parameter degree contributing to the selected exponent. Apply the ordinary analytic cluster transformation rule to the finite analytic families $F(u)$ and $\widetilde Q(u)=\mathsf C_{\mathrm{trunc}}(u)F(u)$. The Taylor coefficients of $\widetilde Q$ and the formal $Q$ agree through that order, so the coefficient identity of residues also agrees. Substituting back the positive Hahn monomials proves \eqref{eq:transform} at the chosen exponent. See \cref{app:transform} for the truncation argument in detail.

\emph{The separated residue is a torus integral.} For the leading tuple $(z_1^\mu,\ldots,z_n^\mu)$, ordinary residues are coefficient extraction. Expanding each $Q_i^{-1}$ around $z_i^{-\mu}$ shows that $\Res_Q$ is the iterated coefficient of $z_1^{-1}\cdots z_n^{-1}$. After the admissible scaling, integration of an ordinary Laurent monomial on each unit circle extracts exactly those same coefficients. All rearrangements are strongly summable by the already verified support bounds. The powers of $r_i$ cancel in a residue term. This proves \eqref{eq:torus}.

The left side is already independent of radii and of $\mathsf C$, and has the support stated in \eqref{eq:AGresidue}. Equality therefore gives all remaining assertions. Independence for any other nonnegative-support representing matrix can also be proved by the same transformation argument; no assertion of uniqueness of that matrix is needed.
\end{proof}

\begin{remark}[What was changed, and what was not]\label{rem:toruslimits}
Formula \eqref{eq:torus} does \emph{not} assert that the original quotient $H/(F_1\cdots F_n)$ is regular on the chosen product torus. A hypersurface $F_i=0$ can meet that torus. We deliberately replaced the denominators by the univariate $Q_i$ and inserted $\det\mathsf C$. Omitting the determinant is generally wrong. Nor have we identified a real $n$-torus with the full boundary of a complex polydisk. The integration cycle and its real dimension are exactly those specified in the theorem.
\end{remark}
Using univariate polynomials in a zero-dimensional ideal to compute multidimensional residues is a classical method \cite[\S0]{CDS}. The additional conclusion here is that this method admits a legitimate support-controlled contour interpretation for the radius-free Hahn algebra. No claim is made that separating denominators or residue transformation are new ideas.

\subsection{Trace, degree, and Stokes on the separating complement}
Combining \eqref{eq:torus} with the supplied trace theorem yields the exact formula
\begin{equation}\label{eq:torustrace}
\frac1{(2\pi i)^n}\HInt_{\mathbb T_r}
 \frac{H J_F\det\mathsf C}{\prod_iQ_i}\,dz_1\wedge\cdots\wedge dz_n
=\sum_a e_aH(a).
\end{equation}
For $H=1$, the right side is the ordinary integer $\mu$. Thus this higher-dimensional contour has a finite local-degree interpretation even when the actual points split across several infinitesimal ranks.

On a common scaled neighborhood avoiding the coordinate divisor $\prod_iQ_i=0$, the displayed meromorphic top form is holomorphic and closed. Its admissible deformations obey \cref{thm:movingStokes}; their periods agree when the intervening chain avoids that divisor. This is a genuine Stokes explanation of contour deformation, in addition to the algebraic radius-independence proof. It is not a statement about arbitrary fine-continuous hypersurfaces in the full class plane.

\section{Rational contour periods at arbitrary scales}\label{sec:rational}
The Hahn-coherent construction relies on an ordinary base or a useful infinitesimal chart. Rational functions admit a different, fully algebraic closed-period theory over any fixed real closed workspace. It does not require an infinite sum or a limiting partition.

\subsection{Algebraic winding}
For an oriented polygonal cycle $C$ in $F^2$ and $a\notin C$, choose a ray from $a$ meeting the edges transversely and avoiding the vertices. Such a generic ray exists because only finitely many directions are excluded. Count its crossings with sign according to the orientation of the edge relative to the ray. With the convention that a counterclockwise circle has index $+1$ at an interior point, this gives
\[
w_F(C,a)\in\Z.
\]
The finite crossing rules show independence of the generic ray: crossings created or destroyed at a tangency occur in oppositely signed pairs, and a passage through a vertex transfers the count between adjacent edges. Definable degree extends this construction to semialgebraic cycles, with the same normalization and homotopy invariance in the punctured field plane \cite{BO}. In particular, a positively oriented semialgebraic Jordan curve has index $1$ on its bounded component and $0$ on its unbounded component.

For polynomial and rational maps on rectangles, Cauchy indices and Sturm chains give a purely real-algebraic argument principle. Eisermann develops the effective program \cite{Eis}; Perrucci--Roy supply corrected product formulae and a systematic treatment of boundary cases \cite{PR}. We use only the boundary-free case. Fractional conventions for zeros on an edge or vertex are not being interpreted as contour integrals through poles.

The word ``effective'' needs a qualification in the present setting. Finite Sturm and sign algorithms can use an ordered-field arithmetic oracle. They do not make arbitrary set-supported Hahn series computable or provide an algorithm for deciding every comparison in an arbitrary presentation of $\Gamma$.

\subsection{Definition and complete rational period theorem}
\begin{definition}[Algebraic rational contour functional]\label{def:algperiod}
Let $R\in K(z)$, and let $C$ be a finite oriented semialgebraic cycle avoiding its poles. Define
\begin{equation}\label{eq:algperiod}
I_C^{\mathrm{alg}}(R\,dz)
=2\pi i\sum_{a\text{ a finite pole of }R}
 w_F(C,a)\res_a(R\,dz).
\end{equation}
The sum is finite. The normalization $2\pi i$ uses the distinguished copy of $\C$ in $K$.
\end{definition}
This is an algebraic definition, not a residue theorem deduced from an unspecified Riemann integral. Over an abstract real closed field without a chosen copy of the real constants, one should instead normalize periods formally or work with the normalized value $I/(2\pi i)$.

\begin{theorem}[Rational Cauchy and argument principles]\label{thm:rational}
The functional \eqref{eq:algperiod} is $K$-linear and additive in cycles. It changes sign under orientation reversal, vanishes on rational exact differentials, and is invariant under definable deformations avoiding the poles. If the indices at all poles vanish, the integral is zero. For a nonzero rational function $R$ with no zero or pole on $C$,
\begin{equation}\label{eq:algargument}
\frac1{2\pi i}I_C^{\mathrm{alg}}\left(\frac{R'}R\,dz\right)
=\sum_a w_F(C,a)\ord_aR,
\end{equation}
where negative orders count poles. On every common admissible ordinary or scaled contour where the Hahn construction and the algebraic winding interpretation both apply, the two rational period functionals agree.
\end{theorem}
\begin{proof}
Linearity and the cycle rules follow from the same properties of residues and winding. The residue of a rational derivative is zero, and winding is unchanged by a pole-avoiding definable homotopy. The logarithmic derivative has residue $\ord_aR$ at each zero or pole, proving \eqref{eq:algargument}.

For agreement, take finite partial fractions of $R$. Polynomial terms and higher-order pole terms have rational primitives, whose admissible Hahn integrals over closed cycles vanish. For a simple-pole term, the Hahn integral is $2\pi i$ times its winding by \cref{thm:winding} and affine rescaling, or by the direct geometric expansion when the pole is outside the normalized finite chart. These winding numbers agree with the algebraic indices on the stated common contours. Summing the finitely many contributions proves agreement.
\end{proof}
For a simple positively oriented Jordan cycle the right side of \eqref{eq:algargument} is the number of inside zeros minus inside poles, with multiplicities. This is an exact arbitrary-scale counterpart of the classical argument principle. It neither supplies open-path logarithm choices nor integrates every nonrational fine-analytic function.

\subsection{Why residues are the complete rational obstruction}
\begin{proposition}[Rational de Rham quotient]\label{prop:rationalcohomology}
Over the algebraically closed characteristic-zero field $K$,
\begin{equation}\label{eq:rationalcohomology}
\frac{K(z)\,dz}{dK(z)}
\simeq\bigoplus_{a\in K}K\left[\frac{dz}{z-a}\right].
\end{equation}
The coordinates of a class are its residues at finite points.
\end{proposition}
\begin{proof}
Partial fractions express every rational differential as a polynomial differential plus finitely many pole terms. Polynomial terms integrate by division by positive integers. For $j\ge2$, $(z-a)^{-j}dz$ is the derivative of $-(z-a)^{-j+1}/(j-1)$. Only the simple poles remain. The residue of a derivative is zero, so a nonzero linear combination of distinct simple-pole terms cannot be exact. This proves both spanning and independence.
\end{proof}
There is no additional independent generator at infinity: its residue is the negative sum of the finite residues. Formula \eqref{eq:algperiod} is therefore the winding pairing with the entire rational de Rham quotient, rather than an arbitrary selection of part of the rational information.

The theorem also explains the difference from the fine-local category. A global fine-analytic primitive of $1/z$ may exist after a phase convention, but no rational primitive exists. The cohomology being computed in \eqref{eq:rationalcohomology} is rational cohomology of a differential algebra, not unrestricted fine-local cohomology.

Finite rational and semialgebraic data over $\SC$ can always be placed in a set-sized real closed workspace. The relevant roots already lie in its algebraic closure $K$, and enlarging the workspace does not change finite sign computations or the residue formula. This supplies an especially robust full-surcomplex interpretation for finite algebraic input.

\section{Worked examples}\label{sec:examples}
\subsection{A Cauchy integral with no common ordinary radius}
Take $\Gamma=\Q$ and the germ \eqref{eq:shrinking}. Set $a=t^2$ and $r=t^{1/2}$. The microscopic Cauchy formula gives
\begin{equation}\label{eq:radiusfreeexample}
\frac1{2\pi i}\HCoint_{|z|=t^{1/2}}
 \frac{\displaystyle\sum_{n\ge1}t^n/(1-nz)}{z-t^2}\,dz
=\sum_{n\ge1}\frac{t^n}{1-nt^2}.
\end{equation}
The right side begins
\[
t+t^2+2t^3+3t^4+5t^5+9t^6+16t^7+31t^8+\cdots.
\]
Its coefficient at exponent $m$ is the finite sum of $n^k$ over positive integers $n$ and nonnegative integers $k$ satisfying $n+2k=m$. The original coefficient functions have poles at $1/n$, so there is no positive ordinary disk carrying them all. After $z=t^{1/2}w$, however, each Hahn coefficient is a polynomial in $w$, and the integral is well defined. The contour construction is doing more than applying Cauchy on a nonexistent shared disk.

\subsection{Two poles in different valuation ranks}
Let $\Gamma=\Q+\Q\omega$, with its inherited surreal order, and put
\[
R(z)=\frac1{z-t}+\frac1{z-t^\omega}.
\]
Both residues are $1$. Three positively oriented microscopic circles give
\begin{equation}\label{eq:rankexample}
\frac1{2\pi i}\HCoint_{|z|=t^s}R(z)\,dz
=\begin{cases}
2,&s=\tfrac12,\\
1,&s=2,\\
0,&s=2\omega.
\end{cases}
\end{equation}
For the middle circle, $t^\omega$ lies inside but $t$ lies outside, because $\omega>2>1$ reverses the magnitude order of the monomials. The last circle is smaller than both pole scales. Each assertion follows either by geometric expansion in the appropriate ratio or by algebraic winding and \cref{thm:rational}. A single real-valued valuation that discards the higher rank cannot retain all of these distinctions.

\subsection{A nonzero infinitesimal area from Stokes}
Let $r\in F_{>0}$ and parameterize the circle by $z=r e^{i\theta}$. For the polynomial real-analytic form $\bar z\,dz$,
\begin{equation}\label{eq:areaexample}
\HCoint_{|z|=r}\bar z\,dz
=2\pi i r^2
=2i\HInt_{r\mathbb D}dx\,dy.
\end{equation}
The area integral is defined by pulling back along the ordinary unit disk under the affine scaling, whose Jacobian is $r^2$. Since $d(\bar z\,dz)=2i\,dx\wedge dy$, this is exactly Stokes. Taking $r=t^s$ gives a nonzero infinitesimal area; taking $r$ infinite gives an infinite value. No measure on all fine-topological subsets is asserted. The same polynomial parameter calculus suffices for every fixed scale in this example.

\subsection{A four-point cluster and its separating torus}\label{subsec:fourpoints}
Consider the two-variable example from \cite[\S11]{AG},
\begin{equation}\label{eq:fourF}
F_1=x^2-Ay,\qquad F_2=y^2-Bx,
\qquad A,B\in\m_K\setminus\{0\}.
\end{equation}
The ordinary leading algebra is $\C\{x,y\}/(x^2,y^2)$ with basis $1,x,y,xy$ and length four. The coordinate characteristic polynomials are
\[
Q_x=x^4-A^2Bx,\qquad Q_y=y^4-AB^2y.
\]
Here an explicit transformation matrix is
\begin{equation}\label{eq:exampleC}
\mathsf C=
\begin{pmatrix}
x^2+Ay&A^2\\ B^2&y^2+Bx
\end{pmatrix},\qquad
\begin{pmatrix}Q_x\\Q_y\end{pmatrix}
=\mathsf C\begin{pmatrix}F_1\\F_2\end{pmatrix}.
\end{equation}
Its determinant is
\begin{equation}\label{eq:exampleD}
D=x^2y^2+Bx^3+Ay^3+ABxy-A^2B^2.
\end{equation}
Thus the radius-free residue functional is realized by
\begin{equation}\label{eq:exampletorus}
\Res_F(H)=\frac1{(2\pi i)^2}\HInt_{\mathbb T_r}
 \frac{H(x,y)D(x,y)}{(x^4-A^2Bx)(y^4-AB^2y)}\,dx\wedge dy.
\end{equation}
Any positive exponents satisfying
\[
3s_x<2v(A)+v(B),\qquad 3s_y<v(A)+2v(B)
\]
provide admissible radii. With $A=t$ and $B=t^\omega$, both $s_x=s_y=1/4$ work.

The roots are $(0,0)$ and
\[
\left(\rho\zeta^j,\frac{\rho^2}{A}\zeta^{2j}\right),
\quad j=0,1,2,\qquad \rho^3=A^2B,
\]
where $\zeta$ is a primitive cube root of unity. The Jacobian is $J_F=4xy-AB$. Its values are $-AB$ at the origin and $3AB$ at the other three points. Consequently the four simple residue weights of $H=1$ are
\[
-\frac1{AB},\quad\frac1{3AB},\quad\frac1{3AB},\quad\frac1{3AB}.
\]
They are individually infinite, but their sum is zero. In contrast,
\begin{equation}\label{eq:examplesres}
\Res_F(xy)=1,\qquad \Res_F(J_F)=4.
\end{equation}
The separating-torus formula reproduces the same cancellation. In the quotient algebra one also finds $D=ABJ_F$; the exact polynomial identity has a remainder in $(F_1,F_2)$.
The separated system $(Q_x,Q_y)$ has $16$ simple grid points, whereas the original system has only four. At each of the other $12$ points, $\mathsf C F=0$ with $F\ne0$, so $\det\mathsf C=0$ and its separated simple residue vanishes. This makes the role of the determinant particularly explicit: changing the denominators without it would count the wrong zero scheme.

For an independent coefficient check, expand
\[
\frac1{Q_x}=x^{-4}\sum_{k\ge0}(A^2B)^k x^{-3k},\qquad
\frac1{Q_y}=y^{-4}\sum_{\ell\ge0}(AB^2)^\ell y^{-3\ell}.
\]
These are admissible on the chosen torus. The integral extracts the coefficient of $x^{-1}y^{-1}$. On the other hand, \eqref{eq:AGresidue} gives, for $p,q\ge0$,
\begin{equation}\label{eq:monomialres}
\Res_F(x^py^q)=
\begin{cases}
A^{(2p+q-3)/3}B^{(p+2q-3)/3},&
\frac{2p+q-3}{3},\frac{p+2q-3}{3}\in\N,\\
0,&\text{otherwise}.
\end{cases}
\end{equation}
The accompanying exact-arithmetic script compares these two procedures on a finite grid of monomials. The general equality is proved by \cref{thm:torus}, not inferred from the tests.

\section{Other integration programs and their actual scope}\label{sec:literature}
The preceding constructions are not the only reasonable responses to the loss of Archimedean topology. Several established theories address nearby problems. Their existence is important, but importing their theorem names without their hypotheses would be misleading.

\subsection{Definable Stokes and integration beyond the ground field}
Ohmoto--Shiota prove a $C^1$ triangulation theorem for semialgebraic sets and formulate integration and Stokes using finite differentiable simplices \cite{OS}. Their Proposition 4.2 is a Stokes formula for the relevant chains, and Theorem 4.1 establishes triangulation independence. This is a close geometric counterpart to the present question.

Their discussion of non-Archimedean real closed fields also warns that integral values need not remain in the given field; a larger target enters. The related measure theory of Ma\v r\'{\i}kov\'{a}--Shiota takes values in an ordered semiring built using a quotient and Dedekind completion \cite{MS}. A semiring-valued measure is not automatically a signed $K$-valued complex line integral. Thus definable Jordan and definable Stokes are meaningful established frameworks, but neither licenses an unqualified claim that every definable form on $F^n$ has a canonical integral in $F$ or $K$.

For the present Hahn-coefficient class, \eqref{eq:Hintegral} and \eqref{eq:Stokes} specify an actual $K$-valued target and avoid that ambiguity. They achieve this by restricting the coefficient structure and parameter chains, rather than by solving the unrestricted definable integration problem.

\subsection{Berkovich spaces and real differential forms}
Chambert-Loir--Ducros construct a differential-form and current theory on Berkovich analytic spaces, with real $(p,q)$-forms and Stokes-type identities \cite{CLD}. For example, Proposition 5.11.1 of the expanded 2025 version gives a tropical Stokes formula at the calibrated tropical level. This is a substantial existing non-Archimedean analogue of generalized Stokes.

Its objects are different from the coefficientwise $K$-valued $dz$ integrals above. Berkovich geometry adds analytic points and uses a real-valued absolute value; its differential-form integrals have a different coefficient and geometric interpretation. For a Hahn field whose value group has higher rank, the full valuation cannot simply be regarded as a real-valued rank-one valuation. Passing to a suitable specialization or coarsening changes what scales are visible. A Berkovich counterpart therefore requires selecting such a setting; it is not a topology-free theorem on the full $\No[i]$.

\subsection{Analytic singular chains and perfectoid integration}
Mihara develops non-Archimedean analytic singular homology using cosimplicial perfectoid spaces and proves a Stokes formula for integrable analytic simplices \cite{Mihara}. In the cited version, Theorem 5.15 relates integration of $d\omega$ to the alternating sum over faces. This supplies a genuine established example in which non-Archimedean integration has its own analytic chain theory rather than ordinary real paths.

A decisive hypothesis is that the valued field has positive residue characteristic. Our field $\C((t^\Gamma))$ has residue characteristic zero. Consequently the cited theorem is not a direct theorem about the present Hahn workspace. Its chain-theoretic architecture is relevant, but replacing its perfectoid and characteristic assumptions would require a new construction and proof.

\subsection{Why a roots-of-unity limit does not automatically transfer}
Shnirelman-type integration and later $p$-adic line-integral theories use arithmetic samples rather than real-parameter contours. Diamond's theory, for example, employs normalized sums over $p$-power roots of unity on $p$-adic arcs and proves Cauchy-type results under its convergence hypotheses \cite{Diamond}.

There is a simple obstruction to copying a familiar sampling prescription into the Hahn valuation topology. On the ordinary unit circle let
\[
f(z)=\frac1{2-z},\qquad
A_N(f)=\frac1N\sum_{\zeta^N=1}\zeta f(\zeta).
\]
Logarithmic differentiation of $X^N-1$ gives
\begin{equation}\label{eq:sampling}
A_N(f)=\frac1{2^N-1}.
\end{equation}
These values tend to zero in the ordinary complex topology, as appropriate for a holomorphic function inside the circle. In a Hahn field with coefficient field $\C$, however, every nonzero value in \eqref{eq:sampling} has valuation zero. The sequence therefore does not converge to zero in its valuation topology, much less in the full fine topology. One could choose an ordinary coefficientwise limit instead, but that is precisely a different definition of convergence. At higher rank, even sampling errors with valuations increasing by a fixed positive increment need not become cofinal.

This example does not refute $p$-adic integration. It pinpoints why its topology, sample geometry, and hypotheses cannot simply be replaced by the word ``surcomplex.''

\subsection{Nonstandard comparisons and resurgent extensions}
Ducros--Hrushovski--Loeser construct a comparison between complex and non-Archimedean integrals using nonstandard models and compatible differential-form structures \cite{DHL}. Their work demonstrates that rigorous bridges between Archimedean and non-Archimedean integration exist, but also illustrates how much additional structure a bridge requires. Real closedness alone is not a transfer principle for arbitrary analytic functions, manifolds, or hyperfinite sums. A hyperfinite construction must be carried out in a specified elementary enlargement with its internal objects.

Costin--Ehrlich develop extension and integration on the surreals for substantial classes including resurgent and related analytic functions \cite{CE}. Such operators are relevant to sectorial extensions and behavior beyond a single common Hahn chart. They do not, by themselves, specify a complex contour atlas with branch transport and monodromy. To use them for the present purpose, one must prove compatibility of the chosen extension, differentiation, and contour or primitive assignments on overlaps.

These comparisons suggest useful directions, but do not replace the concrete support checks made in \cref{sec:deformations,sec:radiusfree,sec:tori}.

\section{What has been established, and what remains to construct}\label{sec:conclusion}
There are three distinct positive answers to the three named classical theorems. Jordan separation survives either after passing to the ordinary standard-part shadow or in semialgebraic/definable topology over a real closed workspace. Cauchy theory survives for Hahn-coherent functions and, after infinitesimal rescaling, for the radius-free Hahn germ algebra. Generalized Stokes survives exactly for Hahn differential forms on ordinary chains and their admissible positive-support deformations.

The separating-torus theorem gives a further bridge rather than merely another analogy. Given the finite-algebra deformation theorem from the supplied geometric manuscript, its residues admit actual coefficientwise contour representatives on microscopic product tori. Characteristic polynomials locate admissible coordinate boundaries, the transformation matrix supplies the necessary determinant, and strong support accounting replaces a nonexistent common ordinary radius. This construction treats singular finite schemes as well as simple roots and keeps every valuation rank.

For rational functions, finite residues paired with algebraic winding supply a closed-period theory at arbitrary scales. This agrees with Hahn integration on their common admissible contours and is complete at the level of rational differentials modulo rational exact forms. It does not determine open-path logarithms or a universal integral for arbitrary fine-local analytic data.

\subsection{Four concrete continuation problems}
A first task is an \emph{admissible scale atlas}: choose a useful class of transition maps, prove strong pullback compatibility, and identify when chartwise contour assignments glue. Equality of actual pulled-back Hahn data is a sufficient criterion here; a flexible global criterion allowing re-expansion is not established.

A second task is a \emph{comparison theorem with definable integration}. One should identify a nontrivial common class of functions and chains for which a selected definable integral and the Hahn integral agree, including their target-field conventions. Agreement of rational periods is a beginning, not a general comparison theorem.

A third task is \emph{sectorial transcendental contour theory}. Essential singularities may require a selected summation or extension operator and compatible monodromy data. For instance, the formal Laurent expansion of $e^{1/z}$ at a nonzero infinitesimal $z$ has descending valuations and is not strongly summable \cite[\S13]{SA}. A chosen extension may exist, but fine analyticity alone does not identify it or its periods.

A fourth task is \emph{higher-rank analytic chains}. Rank-one real-valued non-Archimedean geometries and positive-residue-characteristic perfectoid chains cannot be substituted unchanged for arbitrary-rank complex-residue Hahn workspaces. The finite separating-torus construction provides test cases and exact invariants for any proposed broader chain theory.

These are formulation and compatibility problems left beyond the present results, not claims that named published conjectures have been solved or remain open. The main lesson is precise: the useful analogue is obtained by specifying the geometry and the integral together, not by retaining every classical word while changing only the scalar field.

\appendix
\section{Why finite-parameter testing proves residue transformation}\label{app:transform}
This appendix supplies the technical point in \cref{thm:torus} without assuming an infinite-parameter contour for the original germ family.

\subsection{Finite dependency}
In the supplied normal-form construction, an ordinary linear splitting for $R_n/(f)$ is extended coefficientwise. The perturbed inverse is a Neumann series in operators carrying positive perturbation support. Therefore each coefficient of $M_i$, $Q_i$, and the chosen $C_{ij}$ is built from a finite sum of positive-support words. The residue expansion \eqref{eq:AGresidue} has the same property. For a numerator with support $S$, a coefficient at $\lambda$ only uses decompositions
\[
\lambda=\sigma+e_1+\cdots+e_r,\qquad \sigma\in S,\ e_j\in E.
\]
There are finitely many such words, including their possible lengths. This does \emph{not} say that only finitely many exponents lie below $\lambda$.

Fix a coefficient to be checked in \eqref{eq:transform}. Retain all ordinary source germs and perturbation germs occurring in its finite dependency sets on both sides. Treat numerator source coefficients separately by linearity. Replace the finitely many retained positive perturbation monomials by independent parameters $u_1,\ldots,u_b$. The normal-form formulae now define formal power series in $u$ with ordinary analytic germ coefficients. Their formal identities include
\[
Q(u)=\mathsf C(u)F(u),\qquad Q_i(0,z)=z_i^\mu.
\]
All omitted input terms are irrelevant to the selected coefficient by their construction.

\subsection{Why a formal matrix causes no convergence gap}
The auxiliary ordinary linear splitting need not be continuous in an analytic norm. Thus $\mathsf C(u)$ need not be asserted convergent as an analytic parameter series. Let $N$ exceed every total parameter degree needed for the selected coefficient, and let $\mathsf C_{\le N}(u)$ be its finite truncation. Its finitely many coefficient germs do have a common representative neighborhood. Define the \emph{actual finite analytic family}
\[
\widetilde Q(u,z)=\mathsf C_{\le N}(u,z)F(u,z).
\]
It has ordinary leading tuple $(z_1^\mu,\ldots,z_n^\mu)$ and agrees with the formal $Q(u,z)$ modulo $(u)^{N+1}$. The tuple $F(u,z)$ has leading isolated complete intersection $f$. After choosing ordinary small representatives, both tuples have finite clusters near zero.

The ordinary residue transformation rule, applied to these finite families and their clusters, gives
\[
\res_{F(u)}^{\mathrm{cluster}}(H)
=\res_{\widetilde Q(u)}^{\mathrm{cluster}}
 \bigl(H\det\mathsf C_{\le N}(u)\bigr).
\]
This is the family form of classical residue transformation; it follows from the local transformation law and additivity over the finitely many points \cite{CDS}. Expanding the reciprocal denominators on ordinary protecting residue cycles identifies their Taylor series with the geometric residue expansions. A Taylor coefficient through total parameter degree $N$ uses only denominator and numerator coefficients through that degree. Consequently the last display proves the formal transformation law through degree $N$ for the original $Q$ and $\mathsf C$.

Finally substitute $u_j=t^{e_j}$ and extract the chosen Hahn coefficient. Several parameter monomials may contribute to the same exponent, but only finitely many do. Their finite sum gives the desired coefficient equality. Repeating the argument for each exponent establishes \eqref{eq:transform}. The auxiliary ordinary cycles may change from exponent to exponent; the theorem's final microscopic separating torus is independently constructed by \cref{lem:scaling} and \eqref{eq:torusradius}.

\section{Hypothesis checklist and source map}\label{app:audit}
\begin{center}\small
\begin{longtable}{@{}L{0.28\textwidth}L{0.62\textwidth}@{}}
\toprule
Assertion & Essential condition or provenance\\
\midrule
\endhead
Fine no-path and no-FTC results & Full fine radii; set-sized ordinary parameters; locally constant counterexample. From \cite[\S3]{SA}.\\[0.45em]
Shadow Jordan theorem & Topology consists of standard-part inverse images. Thickened curve, not the ordinary point set. Derived in \cref{thm:shadow}.\\[0.45em]
Definable Jordan theorem & Semialgebraic or appropriately smooth o-minimal embedding and definable connectedness. Established framework \cite{BO}.\\[0.45em]
Coherent Cauchy formula & Common ordinary coefficient domain, uniform Hahn support, point's standard part strictly off the contour. From \cite[\S6]{SA}.\\[0.45em]
Generalized Hahn--Stokes & Ordinary compact parameter chain, smooth globally supported coefficient forms. Derived coefficientwise from classical Stokes.\\[0.45em]
Moving-chain Stokes & Globally positive well-ordered deformation support; analytic target coefficients, or explicit admissible meromorphic inversion. Derived in \cref{thm:movingStokes}.\\[0.45em]
Radius-free contour theory & Germ algebra $\C\{z\}((t^\Gamma))$ and infinitesimal rescaling. Uses translation and division from \cite{AG}.\\[0.45em]
Finite-cluster torus & Supplied finite algebra and support-controlled normal form; classical residue transformation; inequalities \eqref{eq:torusradius}. Derived in \cref{thm:torus}.\\[0.45em]
Rational arbitrary-scale periods & Finite rational poles, semialgebraic cycle avoiding them, specified algebraic winding. Defined by \eqref{eq:algperiod}.\\[0.45em]
Berkovich/perfectoid comparisons & Different geometric points, absolute values, coefficient targets, or residue characteristic. Literature comparisons, not direct transfer theorems.\\
\bottomrule
\end{longtable}
\end{center}
No step uses a proper-class sum. No proof infers arbitrary-rank convergence from $v(\text{error}_n)\ge n\delta$. No separating-torus proof assumes that all original germs have a common ordinary disk. No use of Jordan separation is allowed to turn a fine-local primitive into a coherent primitive without an additional compatibility argument.

\section{Reproducibility}\label{app:verification}
The accompanying \texttt{verify\_examples.py} uses exact SymPy arithmetic to check the transformation matrix \eqref{eq:exampleC}, its determinant, the quotient relation with the Jacobian, and equality of the two residue computations for all monomials $x^py^q$ with $0\le p,q\le10$. It also checks the three sample residues in \eqref{eq:examplesres}, the expansion in \eqref{eq:radiusfreeexample}, and the rational identity underlying \eqref{eq:sampling} for several degrees. The recorded run passes $167$ exact checks. Its output is included in \texttt{verification.txt}.

These checks verify finite illustrative identities. They do not implement a general arbitrary-rank Hahn field, decide arbitrary surreal comparisons, validate every theorem imported from the supplied manuscripts, or constitute a machine-checked proof of the general results. The separation, support, Stokes, and transformation arguments in the text supply the mathematical proofs.

\begin{thebibliography}{99}
\addcontentsline{toc}{section}{References}

\bibitem{SA}
\emph{Surcomplex Analysis: Infinitesimal Calculus, Hahn-Coherent Holomorphy, and Contour Theory}.
Supplied research manuscript, September 21, 2026.
Source file: \texttt{surcomplex\_analysis(3).tex}.
The present article uses its specified coherent framework, not an assertion of universal fine-topological integration.

\bibitem{AG}
\emph{Infinitesimal Analytic Geometry over the Surcomplex Numbers: Radius-Free Hahn Germs, a Nullstellensatz, and Conservation of Singular Zeros}.
Supplied research manuscript, September 21, 2026.
Source file: \texttt{surcomplex\_analytic\_geometry.tex}.
Especially \S\S2--4, 7--10 and Appendix A.

\bibitem{Poonen}
Bjorn Poonen, ``Maximally complete fields,''
\emph{L'Enseignement Math\'ematique} (2) \textbf{39} (1993), 87--106.
\href{https://math.mit.edu/~poonen/papers/amsval.pdf}{Author's full text}.

\bibitem{Neumann}
B. H. Neumann, ``On ordered division rings,''
\emph{Transactions of the American Mathematical Society} \textbf{66} (1949), 202--252.

\bibitem{BO}
Alessandro Berarducci and Margarita Otero,
``Intersection theory for o-minimal manifolds,''
\emph{Annals of Pure and Applied Logic} \textbf{107} (2001), 87--119.
\href{https://doi.org/10.1016/S0168-0072(00)00027-0}{doi:10.1016/S0168-0072(00)00027-0}.
\href{https://arpi.unipi.it/handle/11568/186193}{Institutional record}.

\bibitem{OS}
Toru Ohmoto and Masahiro Shiota,
``$C^1$-triangulations of semialgebraic sets,''
\emph{Journal of Topology} \textbf{10} (2017), 765--775.
\href{https://doi.org/10.1112/topo.12024}{doi:10.1112/topo.12024};
\href{https://arxiv.org/abs/1505.03970}{arXiv:1505.03970}.
Especially \S4, Theorem 4.1, Proposition 4.2, and Remark 4.5.

\bibitem{GH}
Victor Guillemin and Peter J. Haine,
\emph{Differential Forms}, World Scientific, 2019.
\href{https://doi.org/10.1142/11058}{doi:10.1142/11058}.
\href{https://math.mit.edu/classes/18.952/2019SP/18.952_book_2019.pdf}{MIT full-text edition}.
Theorem 4.6.1 gives the ordinary Stokes theorem used here.

\bibitem{CDS}
Eduardo Cattani, Alicia Dickenstein, and Bernd Sturmfels,
``Computing multidimensional residues,'' in
L. Gonz\'alez-Vega and T. Recio (eds.),
\emph{Algorithms in Algebraic Geometry and Applications},
Progress in Mathematics \textbf{143}, Birkh\"auser, 1996, 135--164.
\href{https://arxiv.org/abs/alg-geom/9404011}{arXiv:alg-geom/9404011};
\href{https://doi.org/10.1007/978-3-0348-9104-2_8}{doi:10.1007/978-3-0348-9104-2\_8}.
\S0 includes transformation, coefficient extraction, and univariate separating denominators.

\bibitem{Eis}
Michael Eisermann,
``The Fundamental Theorem of Algebra Made Effective: An Elementary Real-Algebraic Proof via Sturm Chains,''
\emph{American Mathematical Monthly} \textbf{119} (2012), 715--752.
\href{https://doi.org/10.4169/amer.math.monthly.119.09.715}{doi:10.4169/amer.math.monthly.119.09.715};
\href{https://arxiv.org/abs/0808.0097}{arXiv:0808.0097}.
For corrections and boundary subtleties, see \cite{PR}.

\bibitem{PR}
Daniel Perrucci and Marie-Fran\c coise Roy,
``Algebraic Winding Numbers,'' preprint,
\href{https://arxiv.org/abs/2305.08638}{arXiv:2305.08638}, version 3, July 18, 2024.
Theorems 12--13 and the corrected product-formula discussion.

\bibitem{MS}
Jana Ma\v r\'{\i}kov\'{a} and Masahiro Shiota,
``Measuring definable sets in o-minimal fields,''
\emph{Israel Journal of Mathematics} \textbf{209} (2015), 687--714.
\href{https://doi.org/10.1007/s11856-015-1234-0}{doi:10.1007/s11856-015-1234-0};
\href{https://arxiv.org/abs/1402.4787}{arXiv:1402.4787}.

\bibitem{CLD}
Antoine Chambert-Loir and Antoine Ducros,
\emph{Formes diff\'erentielles r\'eelles et courants sur les espaces de Berkovich},
\href{https://arxiv.org/abs/1204.6277}{arXiv:1204.6277},
expanded version 2, July 25, 2025; first version 2012.
In particular, Proposition 5.11.1 is a Stokes formula in the calibrated tropical framework.

\bibitem{Mihara}
Tomoki Mihara,
``Non-Archimedean Analytic Singular Homology Based On Cosimplicial Perfectoid Spaces and Integration along Cycles,''
\emph{The Ramanujan Journal} \textbf{68} (2025), article 98.
\href{https://doi.org/10.1007/s11139-025-01255-8}{doi:10.1007/s11139-025-01255-8};
\href{https://arxiv.org/abs/1211.1422}{arXiv:1211.1422}, version 2, September 16, 2024.
Theorem 5.15 gives the cited Stokes identity; the residue-characteristic hypothesis is essential to the comparison here.

\bibitem{Diamond}
Jack Diamond, ``$p$-adic Line Integrals and Cauchy's Theorems,''
preprint, 2016.
\href{https://arxiv.org/abs/1605.06484}{arXiv:1605.06484}.

\bibitem{DHL}
Antoine Ducros, Ehud Hrushovski, and Fran\c cois Loeser,
``Non-Archimedean integrals as limits of complex integrals,''
\emph{Duke Mathematical Journal} \textbf{172} (2023), 313--386.
\href{https://doi.org/10.1215/00127094-2022-0052}{doi:10.1215/00127094-2022-0052};
\href{https://arxiv.org/abs/1912.09162}{arXiv:1912.09162}.

\bibitem{CE}
Ovidiu Costin and Philip Ehrlich,
``Integration on the surreals,''
\emph{Advances in Mathematics} \textbf{452} (2024), article 109823.
\href{https://doi.org/10.1016/j.aim.2024.109823}{doi:10.1016/j.aim.2024.109823};
\href{https://arxiv.org/abs/2208.14331}{arXiv:2208.14331}.

\end{thebibliography}
\end{document}
